\section{Related Work}

\textbf{Shot boundary detection.} Shot boundary detection has long been an active research area in video processing. Early systems primarily measured frame-to-frame discontinuities using hand-crafted visual features. Methods relying on histogram differences, edge-change ratios, motion compensation, and thresholding-based heuristics are computationally simple. However, they are prone to mistaking lighting changes, camera motion, or object movement for actual shot boundaries \cite{abdulhussain2018review}. Comprehensive survey studies categorize the major challenges in this domain into abrupt cuts, gradual transitions, illumination changes, object/camera motion, and dataset-dependent evaluation protocols \cite{seridi2020survey,kar2024review}.

\textbf{Deep SBD models.} Following the success of deep learning in vision tasks, deep models were introduced to learn temporal patterns that are difficult to encode manually \cite{tran2015c3d,esteve2023shot}. TransNet models SBD as a dense frame-level prediction task over a sequence of low-resolution frames \cite{soucek2019transnet}. TransNetV2 further improves this concept by incorporating dilated 3-D convolutions, frame-similarity features, color histograms, and dual prediction heads for both single-frame and all-transition labels \cite{soucek2020transnetv2}. These models offer a fast and accurate inference pipeline, establishing them as strong baselines for practical SBD applications.

\textbf{AutoShot and short-video SBD.} Most existing datasets, such as BBC, RAI, and ClipShots, mainly represent documentaries, broadcast content, or longer user-generated videos. Recognizing this gap, AutoShot introduces the SHOT dataset specifically to capture short-video characteristics: brief duration, dense editing, vertical layouts, complex gradual transitions, and heavy special effects \cite{zhu2023autoshot}. AutoShot leverages NAS \cite{white2023_nas1000,ren2020_nas_survey,elsken2019nas} to search over approximately 83.9 million candidate architectures composed of 3-D separable convolution blocks and temporal Transformer blocks. While the searched backbone yields notable F1 improvements on SHOT, the computational cost of the search and retraining phases remains prohibitively high.

\textbf{Imbalance, calibration, and temporal smoothing.} Frame-level SBD is severely imbalanced because boundary frames are rare. Focal Loss reduces the contribution of easy examples and is therefore a plausible alternative to BCE \cite{lin2017focal}. SBD decisions are also threshold-sensitive, making score calibration and temporal stability important. Temperature scaling rescales logits without changing their ranking \cite{guo2017calibration}, while temporal smoothing suppresses isolated peaks. We evaluate these mechanisms separately rather than assuming that a more complex training loss is necessary.

